\section{Direct problem}

\subsection{Variational formulation}

The boundary conditions are:
\begin{itemize}
  \item Left and right boundaries ($\Gamma_{\text{side}}$): homogeneous Neumann condition $\partial_n T = 0$ (insulated walls).
  \item Top boundary ($\Gamma_{\text{top}}$): Dirichlet condition $T = T_c = 343\,\mathrm{K}$.
  \item Bottom boundary ($\Gamma_{\text{bottom}}$): Robin condition $\partial_n T + h(T - T_e) = 0$ (door).\
\end{itemize}

The temperature T inside the oven is given by the stationary heat equation:
\begin{equation}\label{base_eq}
  -\mathrm{div}(k\nabla T) = f
\end{equation}

We introduce a lifting function $\tilde{T} \in H^1(\Omega)$ such that $\tilde{T}|_{\Gamma_{\text{top}}} = T_c$. Then, we write $T = \tilde{T} + T_0$, where $T_0$ satisfies the homogeneous Dirichlet condition on $\Gamma_{\text{top}}$.

By multiplying the equation~\eqref{base_eq} by a test function $v \in V = \{v \in H^1(\Omega) \mid v|_{\Gamma_{\text{top}}} = 0\}$ and integrating by parts, we obtain:
\begin{equation}
\int_\Omega k\,\nabla T \cdot \nabla v \dd x - \int_{\partial\Omega} k\,\partial_n T\, v \dd\sigma = \int_\Omega f\,v \dd x.
\end{equation}
By substituting the boundary conditions (zero Neumann on $\Gamma_{\text{side}}$, Robin on $\Gamma_{\text{bottom}}$, and $v=0$ on $\Gamma_{\text{top}}$) and adding the contribution of the lifting function, we arrive at the variational formulation:

\begin{equation}\label{eq:VF}
  \boxed{
  \text{Find } T_0 \in V \text{ such that } a(T_0,v) = l(v), \quad \forall v \in V
  }
\end{equation}where:
\begin{equation}
a(T_0,v) = \int_\Omega k\,\nabla T_0 \cdot \nabla v \dd x + \int_{\Gamma_{\text{bottom}}} k h\,T_0\,v \dd\sigma,
\end{equation}
\begin{equation}
l(v) = \int_\Omega f\,v \dd x + \int_{\Gamma_{\text{bottom}}} k h\,(T_e - \tilde{T})\,v \dd\sigma - \int_\Omega k\,\nabla \tilde{T} \cdot \nabla v \dd x
\end{equation}

\paragraph{Analysis (Lax-Milgram Theorem)}
Let $V = \{v \in H^1(\Omega) \mid v = 0 \text{ on } \Gamma_{\text{top}}\}$ be equipped with the $H^1$-norm $\|v\|_{H^1}^2 = \|v\|_{L^2}^2 + \|\nabla v\|_{L^2}^2$. Recall that the thermal conductivity $k \in L^\infty(\Omega)$ satisfies $1 \le k(\mathbf{x}) \le 10$ a.e. in $\Omega$, and $h = 1$.

To apply the Lax-Milgram theorem, we must verify three hypotheses:

\begin{enumerate}
  \item \textbf{Continuity of the bilinear form $a(\cdot, \cdot)$:}
  For all $u, v \in V$, using the triangle inequality and the upper bound of $k$, we have:
  \begin{equation*}
      |a(u,v)| \le 10 \displaystyle\left\lvert\int_\Omega \nabla u \cdot \nabla v \dd x \right \rvert + 10 \displaystyle\left\lvert\int_{\Gamma_{\text{bottom}}} uv \dd\sigma \right\rvert.
  \end{equation*}
  By applying Cauchy-Schwarz, and using the Trace inequality for the right integral, we obtain:
  \begin{equation*}
      |a(u,v)| \le 10 \|\nabla u\|_{L^2} \|\nabla v\|_{L^2} + 10 C_{\text{tr}}^2 \|u\|_{H^1} \|v\|_{H^1}.
  \end{equation*}
  Since $\|\nabla v\|_{L^2} \le \|v\|_{H^1}$, there exists a constant $M > 0$ such that $|a(u,v)| \le M \|u\|_{H^1} \|v\|_{H^1}$. Thus, $a(\cdot, \cdot)$ is continuous.

  \item \textbf{Coercivity of $a(\cdot, \cdot)$:}
  For all $v \in V$, since $k \ge 1$ and $k h v^2 \ge 0$, we have:
  \begin{equation*}
      a(v,v) = \int_\Omega k |\nabla v|^2 \dd x + \int_{\Gamma_{\text{bottom}}} k h v^2 \dd\sigma \ge \int_\Omega |\nabla v|^2 \dd x = \|\nabla v\|_{L^2}^2.
  \end{equation*}
  Since $\Omega$ is connected and the Dirichlet boundary $\Gamma_{\text{top}}$ has a strictly positive measure, the Poincaré inequality holds on $V$. 
  
  Thus, there exists a constant $C_P > 0$ such that $\|v\|_{L^2(\Omega)} \le C_P \|\nabla v\|_{L^2(\Omega)}$ for all $v \in V$.
  
  Using the definition of the $H^1$-norm ($\|v\|_{H^1}^2 = \|v\|_{L^2}^2 + \|\nabla v\|_{L^2}^2$), this implies that there exists $\alpha > 0$ such that:
  \begin{equation*}
      a(v,v) \ge \|\nabla v\|_{L^2}^2 \ge \alpha \|v\|_{H^1}^2.
  \end{equation*}
  Thus, the bilinear form $a(\cdot, \cdot)$ is coercive on $V$.

  \item \textbf{Continuity of the linear form $l(\cdot)$:}
  For all $v \in V$, we use the triangle inequality and the bounds of $k$:
  \begin{equation*}
      |l(v)| \le \displaystyle\left\lvert \int_\Omega f v \dd x \right \rvert+ 10 \displaystyle\left\lvert \int_{\Gamma_{\text{bottom}}} (T_e - \tilde{T}) v \dd\sigma \right \rvert+ 10 \displaystyle\left\lvert \int_\Omega  \nabla \tilde{T} \cdot \nabla v \dd x \right\rvert.
  \end{equation*}
  Applying the Cauchy-Schwarz inequality and the Trace inequality:
  \begin{equation*}
      |l(v)| \le \|f\|_{L^2} \|v\|_{L^2} + 10 \|T_e - \tilde{T}\|_{L^2(\Gamma_{\text{bottom}})} C_{\text{tr}} \|v\|_{H^1} + 10 \|\nabla \tilde{T}\|_{L^2} \|\nabla v\|_{L^2}.
  \end{equation*}
  Since $\tilde{T}$ is a fixed function in $H^1(\Omega)$ and $f \in L^2(\Omega)$, all terms multiplying the norms of $v$ are bounded constants. Hence, there exists $C > 0$ such that:
  \begin{equation*}
      |l(v)| \le C \|v\|_{H^1}.
  \end{equation*}
  Thus, $l(\cdot)$ is continuous.
\end{enumerate}

\textbf{Conclusion:} Since $V$ is a Hilbert space, $a(\cdot, \cdot)$ is a continuous and coercive bilinear form, and $l(\cdot)$ is a continuous linear form, the Lax-Milgram theorem guarantees the existence and uniqueness of the solution $T_0 \in V$.

\subsection{Finite element approximation and linear system}

We approximate $V$ by a finite-dimensional space $V_h \subset V$ based on a triangular mesh $\mathcal{T}_h$ of $\Omega$ and $\mathbb{P}_1$ finite elements. Let $(\varphi_j)_{j=1}^N$ be the nodal basis of $V_h$.

We write $T_h = \sum_{j=1}^N T_j\,\varphi_j$ and we seek the $T_j$ such that~\eqref{eq:VF} is satisfied for all $v = \varphi_i$. This leads to the linear system:
\begin{equation}
\mathbf{A}\,\mathbf{T} = \mathbf{b},
\end{equation}
where:
\begin{equation}
A_{ij} = \int_\Omega k\,\nabla\varphi_j\cdot\nabla\varphi_i\dd x + \int_{\Gamma_{\text{bottom}}} h\,\varphi_j\,\varphi_i\dd\sigma,
\qquad
b_i = \int_\Omega f\,\varphi_i\dd x + \int_{\Gamma_{\text{bottom}}} h\,T_e\,\varphi_i\dd\sigma.
\end{equation}
The Dirichlet condition $T = T_c$ on $\Gamma_{\text{top}}$ is imposed by modifying the corresponding rows of $\mathbf{A}$ and $\mathbf{b}$ (elimination of Dirichlet degrees of freedom). The matrix $\mathbf{A}$ is symmetric positive definite, sparse, and the system is solved by a direct solver (LU).

\subsection{FreeFem++ implementation and comments}

The implementation uses FreeFem++ with $\mathbb{P}_1$ elements. The geometry is constructed via the \texttt{border} and \texttt{buildmesh} commands, with refinement around the resistors and the zone $S$.

The diffusion coefficient is defined by \texttt{kk = 1 + 9*(region != regS)}, i.e., $k=1$ in $S$ and $k=10$ elsewhere. The Robin condition on the door (label 1) is handled by the term \texttt{int1d(Th,1)(hc*uh*vh - hc*Te*vh)}, and the Dirichlet condition $T=T_c$ on the top (label 3) by \texttt{on(3, uh=Tc)}.

For $\alpha_i = 10^4$ (i=1,\ldots,6), we observe a higher temperature around the resistors and a quasi-uniform distribution in $S$. We can verify that the temperature indeed increases from the door (cooling condition) towards the top (hot wall).

\subsection{Transient problem (heating phase)}

To describe the temperature rise, we add the unsteady term $\rho c\,\partial_t T$, which gives:
\begin{equation}
\rho c\,\partial_t T - \mathrm{div}(k\nabla T) = f \quad \text{in } \Omega \times (0,t_{\max}).
\end{equation}
We approximate the time derivative by an implicit Euler scheme with step $\Delta t$: $\partial_t T^{n+1} \approx (T^{n+1} - T^n)/\Delta t$. The variational formulation at each time step is:
\begin{equation}
\int_\Omega \frac{\rho c}{\Delta t}T^{n+1}v\dd x + \int_\Omega k\nabla T^{n+1}\cdot\nabla v\dd x + \int_{\Gamma_{\text{bottom}}} hT^{n+1}v\dd\sigma
= \int_\Omega \left(\frac{\rho c}{\Delta t}T^n + f\right)v\dd x + \int_{\Gamma_{\text{bottom}}} hT_e\,v\dd\sigma.
\end{equation}
The initial condition is $T^0 = T_e = 293\,\mathrm{K}$ (oven initially at room temperature). The scheme is unconditionally stable (implicit Euler), and the solution converges to the stationary solution as $t\to+\infty$.